\clearpage\begingroup\NoSourceHead\SourceLeaf{542b}
\renewcommand{\thefootnote}{\arabic{footnote})}\setcounter{footnote}{2}
\ArithmeticHead{85}{R}{VEREINZELTE ANWENDUNGEN DER VARIATIONSRECHNUNG.}
\begin{center}\textls[100]{2. Über die elastische Kurve.}\\
\small (Einzelner Zettel, Werke, Bd. XI 1.)\end{center}
\noindent\hspace*{\parindent}\textsc{Daniel Bernoulli} hatte zuerst das Prinzip aufgestellt und \textsc{Euler}\footnote{\textsc{Euler}, \textit{Methodus inveniendi} etc., Lausanne 1744, S. 246: »Cum \textsc{Daniel Bernoulli} mihi indicasset
se universam vim, quae in lamina elastica incurvata insit, una quadam formula quam \textls[120]{vim potentialem}
appellat, complecti posse; hancque expressionem in curva elastica minimam esse oportere; .\ .\ .«.} mitgeteilt, dass die elastische Kurve als Gleichgewichtslage einer elastischen Lamelle die potentielle Energie (»vis potentialis«) der Lamelle, d. h. das Integral
\begin{equation*}\tag{3}
\int\frac{ds}{\rho^2}
\end{equation*}
zu einem Minimum machen muss, wobei $\rho$ den Krümmungsradius, $ds$ das
\clearpage\SourceLeaf{543}
\ArithmeticHead{86}{L}{BOLZA, GAUSS UND DIE VARIATIONSRECHNUNG.}
\noindent Bogenelement der Kurve bedeutet. Hieran anknüpfend hatte dann \textsc{Euler} im
Anhang zu seiner \textit{Methodus inveniendi lineas curvas maximi minimive proprietate
gaudentes} (S. 245—310) seine allgemeinen Methoden auf das Variationsproblem
angewandt: Unter allen (ebenen) Kurven derselben Länge, die nicht nur durch
zwei gegebene Punkte $A$ und $B$ hindurchgehen, sondern auch in diesen
Punkten von gegebenen Geraden berührt werden, diejenige zu bestimmen, die
das Integral (3) zu einem Minimum macht. In rechtwinkligen Koordinaten
$x,y$ mit $x$ als unabhängiger Variabeln nimmt die Aufgabe die Form an: Das
Integral
\[
\int\frac{q^2\,dx}{(1+p^2)^{\frac52}},\quad p=\frac{dy}{dx},\quad q=\frac{d^2y}{dx^2},
\]
zu einem Minimum zu machen mit der isoperimetrischen Nebenbedingung
\[
\int\sqrt{1+p^2}\,dx=S.
\]
\par\noindent\hspace*{\parindent}Da $x$ und $y$ nicht unter dem Integralzeichen vorkommen, so lassen sich
zwei Integrale der \textsc{Euler}schen Differentialgleichung nach einer allgemeinen,
von \textsc{Euler} herrührenden Bemerkung unmittelbar angeben; eine weitere Integration führt \textsc{Euler} auf die Gleichung
\begin{equation*}\tag{4}
\frac{2\sqrt{\alpha\sqrt{1+p^2}+\beta p+\gamma}}{(1+p^2)^{\frac14}}=\beta x-\gamma y+\delta,
\end{equation*}
worin $\alpha$ die isoperimetrische Konstante, $\beta$, $\gamma$ und $\delta$ Integrationskonstanten sind.
Diese Gleichung vereinfacht er dadurch, dass er durch passende Wahl des
Koordinatensystems $\gamma$ und $\delta$ zum Verschwinden bringt. Die Auflösung der
so vereinfachten Gleichung (4) nach $p$ gibt dann das Resultat in der bekannten Form
\[
dy=\frac{(x^2+b)\,dx}{\sqrt{a^4-(x^2+b)^2}}.
\]
\par\noindent\hspace*{\parindent}Hieraus leitet \textsc{Euler} die Gestalt der elastischen Kurve ab, wobei er im
ganzen neun verschiedene Fälle unterscheidet. Lässt man die Grenzfälle bei
Seite, so subsummieren sich alle unter zwei Hauptfälle, die elastische Kurve
mit Wendepunkten und diejenige ohne Wendepunkte.
\par\noindent\hspace*{\parindent}In dem vorliegenden Fragment gibt nun \textsc{Gauss} eine Lösung desselben
Variationsproblems, die an Eleganz wohl nicht zu übertreffen ist. Indem er
\clearpage\SourceLeaf{544}
\ArithmeticHead{87}{R}{VEREINZELTE ANWENDUNGEN DER VARIATIONSRECHNUNG.}
\noindent einerseits die Bogenlänge $s$ als unabhängige Variable und andererseits den
Tangentenwinkel $\varphi$ als unbekannte Funktion einführt\footnote{Denselben Kunstgriff hat M. \textsc{Born} in seiner Dissertation »\textit{Untersuchungen über die Stabilität der
elastischen Linie in Ebene und Raum}« (Göttingen 1906) S. 15, angewandt. Vgl. auch die Behandlung des
Problems bei \textsc{Kneser}, \textit{Lehrbuch der Variationsrechnung} (1900), S. 223—226.}, gelingt es ihm, das
Problem auf ein isoperimetrisches Problem vom einfachsten Typus mit zwei
isoperimetrischen Bedingungen zurückzuführen, bei denen unter dem Integralzeichen nur eine unbekannte Funktion und ihre erste Ableitung vorkommen.
Das gegebene Problem geht nämlich dadurch in die Aufgabe über, $\varphi$ als
Funktion von $s$ so zu bestimmen, dass das Integral
\[
\int_0^S\left(\frac{d\varphi}{ds}\right)^2ds
\]
ein Minimum wird mit den beiden isoperimetrischen Bedingungen
\[
\int_0^S\cos\varphi\,ds=x_1-x_0,\quad\int_0^S\sin\varphi\,ds=y_1-y_0.
\]
und der Anfangsbedingung:
\[
\varphi(0)=\varphi_0,\quad\varphi(S)=\varphi_1,
\]
wobei $x_0,y_0$ und $x_1,y_1$ die Koordinaten der gegebenen Punkte und $\varphi_0$ und $\varphi_1$
die gegebenen Anfangs-, bezw. Endrichtungen bedeuten. Hieraus folgt die
Differentialgleichung
\[
\frac{d^2\varphi}{ds^2}+A\cos\varphi+B\sin\varphi=0,
\]
wobei $A$ und $B$ die beiden isoperimetrischen Konstanten sind. Eine erste
Integration liefert die Gleichung
\[
\frac1\rho=-Ax-By-C,
\]
in der durch passende Wahl des Koordinatensystems bewirkt werden kann, dass
\[
A=0,\quad C=0,\quad B>0.
\]
Eine zweite Integration liefert dann die Gleichung
\begin{equation*}\tag{5}
\cos\varphi=\frac B2\,y^2+D,
\end{equation*}
\clearpage\SourceLeaf{545}
\ArithmeticHead{88}{L}{BOLZA, GAUSS UND DIE VARIATIONSRECHNUNG.}
\noindent die mit der in der oben angegebenen Weise vereinfachten \textsc{Euler}schen Gleichung (4) bei passendem Wechsel der Bezeichnung identisch wird.
\par\noindent\hspace*{\parindent}Hier wäre nun eine Fallunterscheidung nötig gewesen, je nachdem die
Integrationskonstante $D$, die stets $\EqualAboveLess1$ ist, grösser-gleich oder kleiner als $-1$
ist. Indem \textsc{Gauss} $D=\cos2\vartheta$ ansetzt, beschränkt er sich auf den ersten der
beiden Fälle, den er dann auf elliptische Integrale zurückführt. Setzt man
in den \textsc{Gauss}schen Formeln
\[
u=\frac\pi2-\psi,\quad\sin\vartheta=k,
\]
so gehen sie in die übliche Darstellung der Elastica mit Wendepunkten über.
\par\noindent\hspace*{\parindent}Für die Untersuchung des von \textsc{Gauss} bei Seite gelassenen zweiten Falles
$D<-1$ setze man in der Gleichung (5) $\frac\varphi2=\psi$, $\frac2{1-D}=k^2$; dann liefert ein
dem von \textsc{Gauss} für den ersten Fall angewandten analoges Verfahren die üblichen Formeln für die Elastica ohne Wendepunkte.
\clearpage\endgroup
